

The successful chronic implantation of next-gen BCIs is constrained by extreme power efficiency requirements due to patient safety concerns. This has historically been met through highly inflexible, monolithic ASIC hardware. HALO BCIs diverge from the norm by using various function-specific PEs and the RVV extension to introduce significant flexibility without compromising critical thermal limits. However, a highly optimized software compilation pipeline is required to make HALO and its clinical flexibility and capabilities accessible to the broader neuroscientific community. This project directly addresses this need by engineering an end-to-end toolchain that merges a TypeScript-based source-to-source compiler with a hardware-specific vectorization library, enabling the automatic transformation of high-level MATLAB algorithms into power-optimized RISC-V binaries.

Validation and benchmarking demonstrate that the vectorized matrix library successfully runs on Spike and yields near-consistent performance improvements. Analysis of the data has successfully isolated performance improvements driven by updated cross-compiler toolchains, while simultaneously identifying initial test-harness anomalies. The ts-traversal semantic translation layer has also been benchmarked and timed with successful results across various essential operations.

Future research can take multiple directions from here. First, there are unvectorized kernels flagged in the RISCV-Matrix README that require further work. Second, the hand-off to the BCI lab at the Yale School of Medicine will move the pipeline from simulator-only validation onto the physical HALO prototype under neuroscientist-written full workloads, where the 15 mW power envelope and $1 ^\circ$C thermal limit will provide a real-world check on the cycle-count argument made here.

